\scanpage{16}
\uline{ALGEBRE}, \uline{CHAPITRE II} (3e édition)\hfill n\handwritten{°354}

\handwritten{\S} 8. \uline{Modules et anneaux gradués}.

\uline{Dans ce paragraphe}, \uline{on désignera par} $\handwritten{\Delta}$ \uline{un monoïde} (Chap.I, \handwritten{\S} 1,n\handwritten{°3}) \uline{commutatif}, \uline{noté additivement}, \uline{ayant un élément neutre noté} 0.

1. \uline{Groupes commutatifs gradués} ; \uline{anneaux et modules gradués}.

DEFINITION 1.- \uline{Etant donné un groupe commutatif} $G$, \uline{noté additivement}, \uline{on appelle graduation de type} $\handwritten{\Delta}$ \uline{sur} $G$ \uline{une famille} $(G_{\handwritten{\lambda}})_{\handwritten{\lambda}\handwritten{\in}\handwritten{\Delta}}$ \uline{de sous-groupes de} $G$, \uline{dont} $\unsure{G|G'}$ \uline{soit somme directe}. \uline{L'ensemble} $G$, \uline{muni de la structure définie par sa loi de groupe et sa graduation}, \uline{est appelé groupe} (\uline{commutatif}) \uline{gradué de type} $\handwritten{\Delta}$ .

On dit que $\handwritten{\Delta}$ est le \uline{monoïde des degrés} de $G$. On dit qu'un élément $x\handwritten{\in}G$
est \uline{homogène} s'il appartient à un des $G_{\handwritten{\lambda}}$, \uline{homogène de degré} $\handwritten{\lambda}$ si $x\handwritten{\in}G_{\handwritten{\lambda}}$ .
L'élément 0 est donc homogène de tous les degrés ; mais si $x\handwritten{\ne}0$ est homogène
il appartient à un seul des $G_{\handwritten{\lambda}}$ ; l'indice $\handwritten{\lambda}$ tel que $x\handwritten{\in}G$ est alors appelé
\uline{le degré de} $x$ (ou quelquefois \deleted{le} \uline{poids} de $x$) et se note souvent $\deg(x)$.
Tout $y\handwritten{\in}G$ s'écrit d'une seule manière comme somme $\handwritten{\sum}\limits_{\handwritten{\lambda}}y_{\handwritten{\lambda}}$ d'éléments homogènes avec $y_{\handwritten{\lambda}}\handwritten{\in}G_{\handwritten{\lambda}}$ ; on dit que $y_{\handwritten{\lambda}}$ est la \uline{composante homogène de degré}
(ou simplement la \uline{composante de}\quad \uline{degré} $\handwritten{\lambda}$) de $y$.

\uline{Exemples}.- 1) Etant donnés un monoïde commutatif $\handwritten{\Delta}$ quelconque (ayant un
élément neutre) et un groupe commutatif quelconque $G$, on définit une graduation $(G_{\handwritten{\lambda}})_{\handwritten{\lambda}\handwritten{\in}\handwritten{\Delta}}$ sur $G$ en prenant $G_{\unsure{o|0}}=G$ et $G_{\handwritten{\lambda}}=\handwritten{\{}0\handwritten{\}}$ pour $\handwritten{\lambda}\handwritten{\ne}0$ ; cette graduation est dite \uline{triviale}.

2) Soit $\handwritten{\Delta}'$ un second monoïde commutatif ayant un élément neutre, \inserted{\handwritten{et}} soit $\handwritten{\rho}$ un
homomorphisme de $\handwritten{\Delta}$ dans $\handwritten{\Delta}'$ tel que $\handwritten{\rho}(0)=0$. Soit $(G_{\handwritten{\lambda}})_{\handwritten{\lambda}\handwritten{\in}\handwritten{\Delta}}$ une graduation de type $\handwritten{\Delta}$ sur un groupe commutatif $G$ ; pour $\mu\handwritten{\in}\handwritten{\Delta}'$, soit $G'_\mu$ la
somme des $G_{\handwritten{\lambda}}$ tels que $\handwritten{\rho}(\handwritten{\lambda})=\mu$ ; il est clair que $(G'_\mu)_{\mu\handwritten{\in}\handwritten{\Delta}'}$ est une gra-
% archive sheet number: 15
